\documentclass{csc2059practical}

\setpracticalnumber{2}
\setpracticaltitle{Git, Testing, and the CSC2059 Workflow}
\setpracticalsubtitle{Working like a careful C++ developer}
\setpracticalauthor{Dr Giuseppe Trombino, Queen's University Belfast, School of EEECS}
\setpracticaldate{}

\begin{document}

\maketitle
\tableofcontents
\clearpage

\begin{learningoutcomes}
By the end of this practical, you should be able to:
\begin{itemize}
  \item open and work with a Classroom50 repository in VS Code;
  \item run the public test suite using \texttt{make test};
  \item interpret common compiler and test failures;
  \item commit and push your work frequently using Git.
\end{itemize}
\end{learningoutcomes}

\begin{workflowreminder}
For every small change, follow the CSC2059 workflow:
\begin{center}
\textbf{edit $\rightarrow$ make test $\rightarrow$ interpret errors $\rightarrow$ commit $\rightarrow$ push $\rightarrow$ repeat}
\end{center}
Do not wait until the end of the practical to commit. Git is not a ceremonial gong struck only at sunset.
\end{workflowreminder}

\exercise{Clone and Open the Repository}{Estimated time: 10 minutes}

\begin{exerciseobjectives}
  \item Clone your Classroom50 repository.
  \item Open the folder in VS Code.
  \item Confirm that the expected files are present.
\end{exerciseobjectives}

Open a terminal and clone your repository:

\begin{bashcode}
git clone <your-repository-url>
cd <your-repository-name>
code .
\end{bashcode}

\begin{importantnote}
Use VS Code for this module. Do not use Visual Studio. The build and testing workflow for CSC2059 assumes the supplied Makefile and terminal commands.
\end{importantnote}

\exercise{Run the Public Tests}{Estimated time: 15 minutes}

Run the public tests before changing anything:

\begin{bashcode}
make test
\end{bashcode}

A failing test is not a disaster. It is information. Read it carefully before changing code.

\begin{commonerror}
If you see an error such as \texttt{No rule to make target}, check that you are running the command from the root folder of the repository, not from inside a subfolder.
\end{commonerror}

\begin{tipbox}
Run \texttt{pwd} on macOS/Linux or \texttt{cd} on Windows PowerShell to check your current folder.
\end{tipbox}

\exercise{Implement a Small Change}{Estimated time: 25 minutes}

Edit the relevant C++ file and implement one small piece of functionality.

\begin{cppcode}
int countLetters(const std::string& text) {
    int count = 0;
    for (char c : text) {
        if (std::isalpha(static_cast<unsigned char>(c))) {
            ++count;
        }
    }
    return count;
}
\end{cppcode}

After editing, run:

\begin{bashcode}
make test
git status
git add .
git commit -m "Implement letter counting"
git push
\end{bashcode}

\begin{challengebox}
Can you add your own local test case before checking the supplied public tests again?
\end{challengebox}

\begin{reflectionbox}
Briefly answer the following:
\begin{itemize}
  \item What was the first error message you encountered?
  \item What did it actually mean?
  \item What change fixed it?
\end{itemize}
\end{reflectionbox}

\begin{wrapupbox}
You have practised the full CSC2059 development loop: edit, test, diagnose, commit, push, and repeat. This workflow is deliberately simple, but it is also how larger software systems stay under control.
\end{wrapupbox}

\startappendices

\section{Useful Git Commands}

\begin{bashcode}
git status
git add .
git commit -m "Short meaningful message"
git push
git log --oneline
\end{bashcode}

\section{Build Workflow}

Recommended build command:

\begin{bashcode}
latexmk -lualatex -shell-escape example_practical.tex
\end{bashcode}

\end{document}
